% Generated draft: requires mathematical and source review.
\documentclass[11pt]{article}
\usepackage[margin=25mm]{geometry}
\usepackage{fontspec}
\setmainfont{Latin Modern Roman}
\usepackage{amsmath,amssymb,mathtools,unicode-math}
\setmathfont{Latin Modern Math}
\usepackage{enumitem,needspace,xurl,hyperref}
\setlength{\parindent}{0pt}
\setlength{\parskip}{5pt}
\setlength{\emergencystretch}{3em}
\newcommand{\N}{\mathbb{N}}
\newcommand{\Z}{\mathbb{Z}}
\newcommand{\Q}{\mathbb{Q}}
\newcommand{\R}{\mathbb{R}}
\newcommand{\C}{\mathbb{C}}
\newcommand{\F}{\mathbb{F}}
\newcommand{\A}{\mathbb{A}}
\newcommand{\PP}{\mathbb{P}}
\newcommand{\E}{\mathbb{E}}
\newcommand{\eps}{\varepsilon}
\newcommand{\dd}{\,\mathrm{d}}
\begin{document}
\section*{UnsolvedMath ID: 10400212}
\textbf{Source problem code:} catalogue entry
\subsection{Definitions and mathematical statement}
A knotted trivalent graph (KTG), in the usual setting, is a trivalent graph embedded in \(S^3\) with the ribbon/framing structure required by its operations; knot components may occur in the relevant extended category. The excerpt does not specify the elementary KTGs, operation list or presentation relations.

A Vassiliev invariant of order at most \(d\) extends by crossing differences and vanishes on knots with more than \(d\) double points. A universal invariant has a degree-completed diagram-valued expansion \(Z(K)\) whose degree-\(d\) truncation recovers every order-at-most-\(d\) invariant by linear evaluation, with the intended symbol correspondence.

Problem 12.6 asks for a \emph{new proof} of existence of such a universal invariant, using presentations of knots by KTGs and their operations. The proof must explain how the KTG presentations and their relations yield a well-defined knot invariant and establish universality, rather than simply invoke an existing construction. This is a requested proof method, not a claim that existence of universal Vassiliev invariants is itself an unsettled bare existence problem.

The source does not prescribe the coefficient field, framing/orientation category, diagram target, precise KTG operations or any correction/coherence conventions. Strict compatibility with an arbitrarily selected unmodified operation list is not added. The task retains its KTG route without inventing generators, relations or an operation-level invariant that has not been supplied.
\subsection{Short English statement}
Give a new KTG-based proof of existence of a universal finite-type knot invariant. The KTG generators, operations, relations and normalization must be specified; the goal is a new proof through this presentation method.
\subsection{Sources}
\begin{itemize}
\item Local UnsolvedMath registry, permanent ID 10400212, source code catalogue entry; no proposer or proposal year supplied; updated 2026-08-06.
\item Ohtsuki, \emph{Problems on invariants of knots and 3-manifolds} (2002), Problem 12.6, PDF page 162: \url{https://msp.org/gtm/2002/04/gtm-2002-04-024s.pdf}. Registry access date 2026-07-29; PDF font-structure extraction. Provenance distinguishes 231 numbered bold headings. Present primary-literature status and redistribution terms require review.
\item The supplied summary repeats the KTG-based new-proof request without exact further references or a verified resolution.
\end{itemize}
\end{document}
